\section{Comparison and the roles of the three observations}
\label{sec:comparison}
\subsection{Stationary identities and uniform orbit segments}
The relation between stationary entry times and section return times is
classical. Marklof \cite{Marklof} develops this relation through stationary
point processes and Palm distributions, without an ergodicity assumption.
The suspension identity in Section~\ref{sec:flux} is an elementary version
adapted to a finite physical record. It supplies the preparation factor
and the residual-time integral, not the geometry or mass of its active set.
The calculation starts with the stationary action of an actual segment,
rather than with a supplied independent roof sequence.

The use of length as a billiard generating function, and the relation of
action Hessians to linearized dynamics, likewise belong to classical
variational mechanics. Bolotin and Treschev \cite{Hill} treat Hill-type
determinant identities for discrete and continuous Lagrangian systems,
including symplectic twist maps. Our corner-cofactor identity is a finite
Dirichlet calculation of this kind; it is not claimed as a new general
Hill formula. Its role here is to compare the exponentially small physical
flux with its quadratic value \emph{relatively}, uniformly in the number
of reflections. The trace-norm bound in \eqref{eq:v3-logdet} and the summable
endpoint bridge make this possible on one fixed neighborhood.

For fixed $j$, Morse integration alone gives an onset with a
segment-dependent collar and error. Lemma~\ref{lem:g-channels}
verifies the geometric localization at any separated positively curved
periodic configuration. Lemma~\ref{lem:g-relative} gives a common
Dirichlet chart and relative flux with every fixed parameter derivative
uniform in $j$. The two-dimensional radial integral then determines
both the probability and the leading conditional metric. Their different
information content is explicit in Theorem~\ref{thm:v3-fibers} and
Theorem~\ref{thm:v3-tomography}. The finite source hierarchy and
specified record response follow from that same construction. They
are not independent substitutes for a typical-orbit limit theorem.
The circular renewal comparison in Section~\ref{sec:consequences}
retains the change of onset exponent for independent roofs with the
same marginal.

\subsection{Instability from statistical thresholds and from marked lengths}
B\'alint, De Simoi, Kaloshin and Leguil \cite{BDKL} recover curvatures at
period-two points and Lyapunov exponents of periodic orbits from the marked
length spectrum of open dispersing billiards with three strictly convex
obstacles satisfying the non-eclipse condition. Their datum consists of
labelled periodic lengths. Our coefficient sequence comes instead from
full-equilibrium probabilities of maximal collision counts in a periodic
dispersing billiard. Formula~\eqref{eq:ratio} recovers the normal
orbit's multiplier from these statistical thresholds. This is a different
observable, not a claim to stronger inverse-geometric rigidity. The
identification of the multiplier with the Jacobi recurrence is proved
explicitly in Section~\ref{sec:consequences}.


\subsection{Periodic rare events and physical collision onsets}
Carney, Nicol and Zhang \cite[Theorems 1--2]{CNZ} prove extreme-value
and compound-Poisson limits for visits to shrinking neighborhoods of
regular periodic points, including finite- and infinite-horizon Sinai
billiard maps. Their metric is adapted to stable and unstable foliations;
the extremal index is $1-|DT^p|_{E^u}|^{-1}$, and the associated
cluster multiplicities are geometric. Their observation length diverges
while the target shrinks so that the expected exceedance count stays
finite. Thus statistical detection of periodic instability is already
present in that work; it is not an originality claim of this article.

Our observable is the maximal number of physical impacts in a window
near $jg_e$, under full flow preparation. The collar, relative error and
shape/source derivatives are uniform in $j$, rather than derived from
a rare-visit limit. Theorems~\ref{thm:g-stability} and
\ref{thm:g-marked} verify those claims under noncircular deformations
and at competing threshold surfaces. Neither a rare-visit law nor a
fixed-segment Morse asymptotic supplies this uniform physical statement.
The proof does not invoke a mixing rate. Conversely it does not prove
a compound-Poisson law for long observation sequences.

\subsection{Scalar amplitudes, Prony recovery and conditional metrics}
After the odd-index M\"obius transform, the exact recovery in
Theorem~\ref{thm:g-inverse} is a finite-exponential-sum problem.
Its Hankel and Vandermonde steps are classical Prony machinery.
Batenkov and Yomdin \cite[Theorem 4.5]{BYProny} quantify local
conditioning of confluent Prony systems at regular points, with
bounds depending on node separation and the relevant coefficients.
Our proof retains the elementary recovery argument. It attributes
the geometric work to the amplitude formula and specified transform,
not to a new Hankel algorithm. Proposition~\ref{prop:v3-area-coalescence}
shows directly within the equal-gap family why a finite unlabelled
amplitude vector loses first-order curvature information at the circle.

Theorem~\ref{thm:v3-tomography} does not regularize that same vector.
For a supplied quadratic sublevel the covariance identity is elementary;
the point here is its realization by the physical conditional record,
with a differentiated remainder uniform in the flight number.
The theorem uses two centered Euclidean endpoint variances from
physically selected records. Their limits are the diagonal entries of the inverse
endpoint Hessian divided by three. An odd collision length places them
at different ends of the channel. Their product and ratio determine
both curvatures, with uniform Lipschitz bounds even when the curvatures
coincide and without prior knowledge of the preparation area. This
remains valid with unknown higher boundary jets in the admitted smooth
family. Proposition~\ref{prop:v3-sampling} states the finite-offset
bias, sampling accuracy and exponentially growing preparation cost.
The scalar fiber in Theorem~\ref{thm:v3-fibers} is realized by actual
analytic obstacles with fixed area and gap; it is not a formal change
of two otherwise unverified contact jets.

For comparison in the periodic inverse setting, Finamore and Leguil
\cite[Theorem A]{FinamoreLeguil} prove isometry of diffeomorphic
finite-horizon Sinai tables with the same enriched marked length
spectrum. Their enriched datum is defined using limiting marked lengths
of approximating geodesic flows and includes data not represented by
ordinary periodic billiard trajectories. Its global rigidity conclusion
is different from local contact recovery by conditional variances.
Our theorem allows infinite horizon but does not recover the entire
boundary or assert rigidity from a length spectrum. Neither datum is
silently supplied by the other observation. The marked-length results
of \cite{BDKL} remain a distinct antecedent for curvature recovery in
the open three-obstacle setting.

In the circular specialization the threshold spacing already fixes
$R$ and the multiplier. Corollary~\ref{cor:multiplier} is retained
as an exact amplitude identity, not as an additional independent shape
invariant. The equal-gap theorem instead varies three contact curvatures
while keeping all threshold locations fixed; its area is the symmetric
function proved in \eqref{eq:v3-area}. The general unequal-curvature
count coefficient keeps only $c_0c_1$. Channel selection and endpoint
positions in the marked theorem supply precisely the information
missing from that scalar coefficient, without asserting that a
nonminimal channel remains visible in every ground-onset window.

\subsection{Nonlinear boundary laws and higher jets}
Proposition~\ref{prop:v4-saturation} identifies the exact finite-dimensional
factorization of the leading selected-channel hierarchy. The nonlinear
conclusions do not contradict it. Theorem~\ref{thm:v4-factorization}
passes from a finite relative determinant to two half-line determinants,
with quantitative control of their interaction. Theorem~\ref{thm:v4-law}
then gives a fixed-positive-offset physical probability law and
Corollary~\ref{cor:v4-poles} its first generating-function singularities.
The analytic family in Theorem~\ref{thm:v4-jet-fiber} keeps even the
full leading endpoint metric fixed, while its nonlinear coefficient
changes with derivative $\sqrt3/2$. These are consequences of the
uniform nonlinear construction, not additional copies of the leading
two-curvature inverse.

Higher-order information in billiard data is not in itself a new
inverse-theoretic principle. Osterman \cite{Osterman} considers three
analytic scatterers under a non-eclipse condition. Marked periodic
lengths determine analytic conjugacy near a homoclinic orbit, and,
when two scatterers are supplied, determine the third. That work uses
Birkhoff normal forms and asymptotic expansions to recover analytic
dynamical information. Our data are instead physical probabilities
and conditional endpoint records in a periodic table. The explicit
fourth-jet separation is not a claim to his global recovery conclusion
or to the first appearance of nonlinear jets in an inverse billiard
problem. The relative determinant and full-phase residual-time
integration are indispensable to the particular statistical law here.

There are also two quantitative additions to the leading observations.
Theorem~\ref{thm:v4-radial-stability} measures dispersion by a strictly
convex three-amplitude statistic and proves a positive one-half-H\"older
estimate relative to the circular table. Theorem~\ref{thm:v4-sequence-stability}
shows that this reference-point exponent is sharp even for the complete
sequence in the specified absolute-error topology. Neither is a
uniformly regular Prony inverse at arbitrary coalescences.
Theorem~\ref{thm:v4-acquisition} uses fixed-order polynomial
extrapolation of the uniform physical variance expansion. Its explicit
cost improves the sufficient accuracy exponent of the single-offset
estimator but retains the exponential rarity factor. The calibration
bound in Proposition~\ref{prop:v4-timing} charges the physical-window
error separately. These are quantitative statements for their stated
observations, not claims that additional collision order makes the
leading parameters easier to estimate.

\subsection{Perturbative limit laws and the original Fourier programme}
Demers and Zhang \cite{DemersZhang} construct a functional-analytic
framework for billiard-map perturbations, including scatterer motions and
deformations. Their perturbation estimates and continuity of spectral
data address singularities on a much larger dynamical phase space than
our onset charts. Theorem~\ref{thm:g-marked} does not replace those operator
estimates. It gives every fixed shape derivative of a specified extremal
record, including its threshold singularity, in a regime selected by an
exact physical event. It is not an assertion of all-order differentiability
of the full transfer operator under arbitrary deformations.

Demers, Melbourne and Nicol \cite{DMN} prove martingale approximation and
statistical limit laws for H\"older observables of the time-one map of a
finite-horizon planar periodic Lorentz gas. Those typical-orbit limit laws
and the present maximal-count asymptotics concern different regimes.
In particular, neither their stated central limit and invariance principles
nor our threshold series is a source-differentiated mixed local Edgeworth
expansion for general collision records.

The earlier arithmetic and Fourier-inversion chapter \cite{A2R33} proves
an algebraic separation estimate and integrates four explicitly assumed
frequency envelopes. A mechanical application requires the characteristic
family itself to satisfy those envelopes, together with a central
integrable remainder and the required lattice and covariance conditions.
The present results provide a uniform, source-dependent calculation on a
nontrivial family of actual multi-impact events. They do not estimate the
complementary itineraries or construct the stable-curve operators needed
for that original long-time application. In particular, the sum in
\eqref{eq:summable} ranges over distinct onset windows; it cannot be
substituted for the characteristic function of one unrestricted long
record. Establishing that substitution's actual prerequisites is a
separate mathematical problem, not a consequence of the notation.

Thus the structural conclusion is a collision-order-uniform geometric
threshold law and its nonlinear two-boundary limit, on arbitrary
separated positive-curvature periodic configurations. The scalar
leading amplitude, the leading endpoint metric and the nonlinear law
have successively different information content, separated by explicit
geometric families and quantitative estimates. The
fixed-window theorem in \cite{TC}, the conditional Fourier compiler in
\cite{A2R33}, and this uniform extremal theorem have separate, explicit
hypotheses and conclusions. No implication between them is used without
proof.
